\lecture{7}{15. September 2025}{Euler and Bernoulli equation}

\section{Incompressible Inviscid Flow}
The Navier-Stokes equations introduced in \autoref{sec:navsto} describe fluid behaviour in a wide range of problems. However these are not analytically solvable except in simple cases. Instead of the Navier-Stokes equations we can instead use Euler's equation, which applies to an inviscid ($\mu = 0$) fluid (which itself is an idealized notion as no flow is truly inviscid).

\subsection{Momentum Equation for Frictionless Flow: Euler's Equation}
Euler's equation can be obtained from the Navier-Stokes equations simply by neglecting the viscous term. I.e.
\begin{equation} \label{eq:eul}
  \rho \frac{\mathrm{D}\textbf{V}}{\mathrm{D}t} = \rho \textbf{g} - \nabla p
\end{equation}
This states that for an inviscid flow, the change in momentum of a fluid particle is caused by the body force and the net pressure force. In rectangular coordinates this is
\begin{align*}
  \rho \left( \frac{\partial u}{\partial t} + u\frac{\partial u}{\partial x} + v \frac{\partial u}{\partial y} + w \frac{\partial u}{\partial z}  \right) &= \rho g_x - \frac{\partial p}{\partial x}  \\
  \rho \left( \frac{\partial v}{\partial t} + u\frac{\partial v}{\partial x} + v \frac{\partial v}{\partial y} + w \frac{\partial v}{\partial z}  \right) &= \rho g_y - \frac{\partial p}{\partial y}  \\
  \rho \left( \frac{\partial w}{\partial t} + u\frac{\partial w}{\partial x} + v \frac{\partial w}{\partial y} + w \frac{\partial w}{\partial z}  \right) &= \rho g_z - \frac{\partial p}{\partial w}
.\end{align*}
If the $z$-axis is vertical, then $\textbf{g} = - g \hat{k} \implies g_x = g_y = 0, g_z = -g$.

In general, \autoref{eq:eul} is applicable for any problem in which there is no viscous stresses. This also includes rigid-body motion problems. 


\begin{figure} [ht]
  \centering
  \includegraphics[width=0.5\linewidth]{./figures/f6_1.png}
  \caption{Fluid particle moving along a streamline.}
  \label{fig:f6_1}
\end{figure}

We consider the flow in the $yz$-plane shown on \autoref{fig:f6_1}. Here our goal is to write the equations of motion in terms of the distance along the streamline $s$ and the distance normal to the streamline $n$. The pressure at the center of the fluid element is assumed to be $p$. If we apply Newton's second law in the direction $s$ of the streamline, to the fluid element on \autoref{fig:f6_1} with volume $\mathrm{d}s \, \mathrm{d}n \, \mathrm{d}x$, neglecting viscous forces we obtain
\[ 
\left( p - \frac{\partial p}{\partial s} \frac{\mathrm{d}s}{2}  \right) \, \mathrm{d}n \, \mathrm{d}x - \left( p + \frac{\partial p}{\partial s} \frac{\mathrm{d}s}{2} \right) \, \mathrm{d}n \, \mathrm{d}x - \rho g \sin \beta \, \mathrm{d}s \, \mathrm{d}n \, \mathrm{d}x = \rho a_s \, \mathrm{d}s \, \mathrm{d}n \, \mathrm{d}x
\]
where $\beta$ is the angle between the tangent to the streamline and the horizontal, and $a_s$ is the acceleration of the fluid particle along the streamline. By simplification, we obtain
\[ 
- \frac{\partial p}{\partial s} - \rho g \sin \beta = \rho a_s
.\]
Since $\sin \beta = \frac{\partial z}{\partial s} $, this can be written as
\[ 
- \frac{1}{\rho} \frac{\partial p}{\partial s} - g \frac{\partial z}{\partial s}  = a_s
\]
which holds along any streamline $V = V(s,t)$. We remember that the total acceleration of a fluid particle in the streamwise direction is
\[ 
a_s = \frac{\mathrm{D}V}{\mathrm{D}t} = \frac{\partial V}{\partial t} + V \frac{\partial V}{\partial s}
.\]
Hence Euler's equation in the streamwise direction with the $z$-axis directed vertically upward is
\[ 
- \frac{1}{\rho} \frac{\partial p}{\partial s} - g \frac{\partial z}{\partial s} = \frac{\partial V}{\partial t} + V \frac{\partial V}{\partial s} 
.\]
For steady flow this reduces to
\begin{equation} \label{eq:euleqstrste}
\frac{1}{\rho} \frac{\partial p}{\partial s} = - g \frac{\partial z}{\partial s} - V \frac{\partial V}{\partial s}   
\end{equation}

This indicates that the pressure along a streamline is influenced by the gravitational field and the velocity. The gravitational effect has already been described in \autoref{sec:flusta} -- When $V = 0$, the pressure increases directly proportionally with the change in elevation. In the $xy$-plane, in which there is no influence of gravity \autoref{eq:euleqstrste} reduces to
\[ 
\frac{1}{p} \frac{\partial p}{\partial s}  = - V \frac{\partial V}{\partial s} 
\]
which states that for an incompressible, inviscid flow a decrease in pressure will cause an increased velocity and v.v. 

To obtain Euler's equation in a direction normal to the streamlines, we apply Newton's second law in the $n$ direction to the fluid element. Again, neglecting viscous forces, we get
\[ 
\left( p - \frac{\partial p}{\partial n} \frac{\mathrm{d}n}{2} \right) \, \mathrm{d}s \, \mathrm{d}x  - \left( p + \frac{\partial p}{\partial n} \frac{\mathrm{d}n}{2} \right) \, \mathrm{d}s \, \mathrm{d}x - \rho g \cos \beta \, \mathrm{d}n \, \mathrm{d}x \, \mathrm{d}s = \rho a_n \, \mathrm{d}n \, \mathrm{d}x \, \mathrm{d}s
\]
where $\beta$ is the angle between the $n$-direction and the vertical and $a_n$ is the acceleration of the fluid particle in the $n$-direction
On simplification this becomes
\[ 
- \frac{\partial p}{\partial n} - \rho g \cos \beta = \rho a_n 
.\]
And since $\cos \beta = \frac{\partial z}{\partial n} $ this can be written as
\[ 
- \frac{1}{\rho} \frac{\partial p}{\partial n} - g \frac{\partial z}{\partial n} = a_n
.\]
The normal acceleration of the fluid element is toward the center of curvature of the streamline -- in the $-n$-direction. Hence
\[ 
a_n = - \frac{V^2}{R}
.\]
Therefore, for steady flow, where $R$ is the radius of curvature of the streamline at the point chosen, Euler's equation normal to the streamline is
\[ 
\frac{1}{\rho} \frac{\partial p}{\partial n} + g\frac{\partial z}{\partial n} = \frac{V^2}{R}
.\]
In a horizontal plane this becomes
\[ 
\frac{1}{\rho} \frac{\partial p}{\partial n} = \frac{V^2}{R}
.\]


\subsection{Bernoulli Equation: Integration of Euler's Equation Along a Streamline for Steady Flow}
Compared to the Navier-Stokes equations the Euler's equation for incompressible inviscid flow is mathematically simpler -- however solving it still presents difficulty in complex cases. One convenient approach to circumvent this for a steady flow is to integrate Euler's equation along a streamline.

\subsubsection{Derivation of the Bernoulli Equation Using Streamline Coordinates} \label{sec:bernstre}
Euler's equation for steady flow along a streamline is given in \autoref{eq:euleqstrste} as
\[ 
- \frac{1}{\rho} \frac{\partial p}{\partial s} - g \frac{\partial z}{\partial s} = V \frac{\partial V}{\partial s}
.\]
If a fluid particle moves an infinitesimal distance $\mathrm{d}s$ along a streamline, then
\begin{align*}
  \frac{\partial p}{\partial s} \, \mathrm{d}s &= \mathrm{d}p & &(\text{The change in pressure along } s) \\
  \frac{\partial z}{\partial s} \, \mathrm{d}s &= \mathrm{d}z & &(\text{The change in elevation along } s) \\
  \frac{\partial V}{\partial s} \, \mathrm{d}V &= \mathrm{d}p & &(\text{The change in speed along } s)
.\end{align*}
Thus, by multiplying \autoref{eq:euleqstrste} be $\mathrm{d}s$ we get
\[ 
- \frac{\mathrm{d}p}{\rho} - g \, \mathrm{d}z = V \, \mathrm{d}V \implies \frac{\mathrm{d}p}{\rho} + V \, \mathrm{d}V + g \, \mathrm{d}z = 0 \qquad (\text{along } s)
.\]
Integration of this equation gives
\[ 
\int \frac{\mathrm{d}p}{\rho} + \frac{V^2}{2} + gz = \mathrm{constant} \qquad (\text{along } s)
.\]
To use this one must first find the inverse change in density with pressure $\frac{1}{\rho} \, \mathrm{d}p$. For incompressible flow $\rho = \mathrm{constant}$ so the above equation becomes the Bernoulli equation,
\begin{equation} \label{eq:berneq}
  \frac{p}{\rho} + \frac{V^2}{2} + gz = \mathrm{constant}
\end{equation}
For \autoref{eq:berneq} to be applicable the flow must be steady, incompressible, frictionless and happen along a streamline. 


\subsubsection{Derivation of the Bernoulli Equation Using Rectangular Coordinates}
Instead of using Euler's equation on streamline-coordinate form as in \autoref{sec:bernstre} one could also integrate the vector form of Euler's equation from \autoref{eq:eul} along a streamline. As we are simply interested in the Bernoulli equation we will only consider the case of steady flow.

For steady flow, Euler's equation in rectangular coordinates is
\begin{equation}\label{eq:eulrec}
  \frac{\mathrm{D}\textbf{V}}{\mathrm{D}t} = u \frac{\partial \textbf{V}}{\partial x} + v \frac{\partial \textbf{V}}{\partial y} + w \frac{\partial \textbf{V}}{\partial z} = \left( \textbf{V} \cdot \nabla  \right)\textbf{V} = - \frac{1}{\rho} \nabla p - g \hat{k}
\end{equation}
For steady flow the velocity field is $\textbf{V} = \textbf{V}\left( x,y,z \right)$. The motion of a particle along a streamline is governed by \autoref{eq:eulrec}. During the time interval $\mathrm{d}t$ the particle has a displacement $\mathrm{d}\textbf{s}$ along the streamline.

If we take the dot product of \autoref{eq:eulrec} and the displacement $\mathrm{d}\textbf{s}$ along the streamline, a scalar equation relating pressure, speed, and elevation is obtained along the streamline. I.e.
\[ 
\left( \textbf{V} \cdot \nabla  \right) \textbf{V} \cdot \mathrm{d}\textbf{s} = - \frac{1}{\rho} \nabla p \cdot \mathrm{d}\textbf{s} - g \hat{k} \cdot \mathrm{d}\textbf{s}
\]
where
\[ 
\mathrm{d}\textbf{s} = \mathrm{d}x \hat{i} + \mathrm{d}y \hat{j} + \mathrm{d} z \hat{k} \qquad (\text{along } s)
.\]
Now we can evaluate each of the three terms one-by-one, starting with $-\frac{1}{\rho} \nabla p \cdot \mathrm{d}\textbf{s}$,
\begin{align*}
  - \frac{1}{\rho} \nabla p \cdot \mathrm{d}\textbf{s} &= - \frac{1}{\rho} \left( \hat{i} \frac{\partial p}{\partial x} + \hat{j} \frac{\partial p}{\partial y} + \hat{k} \frac{\partial p}{\partial z}  \right) \cdot \left( \mathrm{d}x \hat{i} + \mathrm{d} y \hat{j} + \mathrm{d}z \hat{k} \right) & &(\text{along } s) \\
   &= - \frac{1}{\rho} \left( \frac{\partial p}{\partial x} \, \mathrm{d}x + \frac{\partial p}{\partial y} \, \mathrm{d}y + \frac{\partial p}{\partial z} \, \mathrm{d}z \right) & &(\text{along } s) \\
  &= - \frac{1}{\rho} \, \mathrm{d}p & &(\text{along } s)
.\end{align*}
Moving on to $- g \hat{k} \cdot \mathrm{d}\textbf{s}$, we get
\begin{align*}
  - g \hat{k} \cdot \mathrm{d}\textbf{s} &= - g \hat{k} \cdot \left( \mathrm{d}x \hat{i} + \mathrm{d}y \hat{j} + \mathrm{d}z \hat{k} \right) &&(\text{along }s)\\
    &= - g \, \mathrm{d}z & &(\text{along }s)
.\end{align*}
Now we just need to evaluate $\left( \textbf{V}\cdot \nabla  \right)\textbf{V}\cdot \mathrm{d}\textbf{s}$. A common vector identity allows us to rewrite this as
\begin{align*}
  \left( \textbf{V}\cdot \nabla  \right)\textbf{V}\cdot \mathrm{d}\textbf{s} &= \left( \frac{1}{2} \nabla \left( \textbf{V} \cdot \textbf{V} \right) - \textbf{V} \times \left( \nabla \times \textbf{V} \right) \right) \cdot \mathrm{d}\textbf{s} \\
      &= \left( \frac{1}{2}\nabla  \left( \textbf{V} \cdot \textbf{V} \right) \right)  \cdot \mathrm{d}\textbf{s} - \left( \textbf{V} \times \left( \nabla \times \textbf{V} \right) \right) \cdot \mathrm{d}\textbf{s}
.\end{align*}
The last term in this is zero, since $\textbf{V}$ is parallel to $\mathrm{d}\textbf{s}$. Hence,
\begin{align*}
  \left( \textbf{V} \cdot \nabla  \right) \textbf{V} \cdot \mathrm{d} \textbf{s} &= \frac{1}{2} \nabla  \left( \textbf{V} \cdot \textbf{V} \right) \cdot \mathrm{d}\textbf{s} &&\left( \text{along }s \right) \\
  &= \frac{1}{2} \nabla \left( V^2 \right) \cdot \mathrm{d}\textbf{s} &&(\text{along }s) \\
  &= \frac{1}{2} \left( \hat{i} \frac{\partial V^2}{\partial x} + \hat{j} \frac{\partial V^2}{\partial y} + \hat{k} \frac{\partial V^2}{\partial z}  \right) \cdot \left( \mathrm{d}x \hat{i} + \mathrm{d}y \hat{j} + \mathrm{d}z \hat{k} \right) &&(\text{along }s)\\
  &= \frac{1}{2} \left( \frac{\partial V^2}{\partial x} \, \mathrm{d}x + \frac{\partial V^2}{\partial y} \, \mathrm{d}y + \frac{\partial V^2}{\partial z} \, \mathrm{d}z \right) &&(\text{along }s)\\
  &= \frac{1}{2} d \left( V^2 \right) &&(\text{along }s)
.\end{align*}
Substituting these three evaluated terms into the expression yields
\[ 
\frac{\mathrm{d}p}{\rho} + \frac{1}{2}d \left( V^2 \right) + g \mathrm{d}z = 0 \qquad (\text{along }s)
.\]
Integrating this we obtain
\[ 
\int \frac{\mathrm{d}p}{\rho} + \frac{V^2}{2} + gz = \mathrm{constant}
.\]
And for constant density this becomes the Bernoulli equation, \autoref{eq:berneq},
\[ 
\frac{p}{\rho} + \frac{V^2}{2} + gz = \mathrm{constant}
.\]
Once again, for \autoref{eq:berneq} to be applicable the flow must be steady, incompressible, frictionless and happen along a streamline. 

\subsubsection{Static, Stagnation, and Dynamic Pressures}
The pressure $p$ which is used in the Bernoulli equation, \autoref{eq:berneq} is the thermodynamic pressure, also called the \textit{static pressure}. This is the pressure measured by an observer riding along with the fluid. 

The \textit{stagnation pressure} is obtained when a flowing fluid is decelerated to zero speed by a frictionless process. For incompressible flow, the Bernoulli equation can be used to relate changes in speed and pressure along a streamline for such a process. Neglecting elevation differences \autoref{eq:berneq} becomes
\[ 
\frac{p}{\rho} + \frac{V^2}{2} = \mathrm{constant}
.\]
The stagnation pressure $p_0$ can then be computed by setting the stagnation speed $V_0 = 0$, as
\begin{equation}\label{eq:stagpre}
  \frac{p_0}{\rho} + 0 = \frac{p}{\rho} + \frac{V^2}{2} \implies p_0 = p + \frac{1}{2}\rho V^2
\end{equation}
This is the mathematical definition of stagnation pressure for incompressible flow. The term $\frac{1}{2}\rho V^2$ is called the \textit{dynamic pressure}. Thus, the stagnation  (or \textit{total}) pressure is the sum of the static pressure and the dynamic pressure.

Solving \autoref{eq:stagpre} for the local flow speed $V$ yields
\[ 
V = \sqrt{\frac{2 \left( p_0 - p \right)}{\rho}}
.\]
Thus if one knows the stagnation pressure and the local pressure one could calculate the corresponding flow speed.
